\subsection{Biological case studies and interpretability}
\label{sec:case_studies}

To demonstrate the interpretability of SSV features and to illustrate representative benchmark behavior, we present three detailed case studies spanning successful prediction, size-selective binding, and failure mode diagnosis.

\textbf{Case 1: AB\_CB---Terminal-rich glycan recognition.} This antibody (5 positive glycans) achieves strong SSV-based prediction performance (MRR = 0.396, Recall@5 = 0.60). SSV preference analysis (Table~\ref{tab:e7_preferences}) identifies a significant preference for glycans with high terminal residue counts (terminal\_proxy: Cliff's $\delta = +0.79$, Cohen's $d = +1.32$, $p_{\text{Holm}} = 0.008$). The five known positive glycans (G00066MO, G17480GI, G66163TI, G74353FF, G99315PE) exhibit mean terminal\_proxy = 3.0 versus background mean = 1.86, reflecting a large effect size that persists under size-matched control (E7b). This preference is biologically plausible: terminal residues (non-reducing ends) are primary epitope sites for antibody recognition, as they present exposed monosaccharide determinants unobstructed by glycosidic linkages. The strong prototype-based ranking demonstrates that a single geometric property---terminal residue abundance---can serve as a discriminative structural fingerprint when binding specificity is narrow.

\textbf{Case 2: AB\_1E3---Size-selective binding to compact glycans.} This antibody (5 positive glycans: G00068MO, G35778AN, G40696TO, G61459JW, G94435QH) shows weaker but significant SSV-based prediction (MRR = 0.156, Recall@5 = 0.40). Preference profiling reveals three correlated size/geometry dimensions: n\_atoms ($\delta = -0.72$), radius\_of\_gyration ($\delta = -0.72$), and max\_pair\_distance ($\delta = -0.74$), all with large negative effects (Cohen's $d \approx -1.0$, $p_{\text{Holm}} < 0.05$). Positive glycans are smaller (mean 3.8 residues vs background 5.6) and more compact (mean compactness 0.337 vs 0.317). Despite strong effect sizes, ranking performance is modest (MRR = 0.156), suggesting that size preference alone is insufficient for accurate prediction in the current benchmark. This may reflect: (1) sequence-level epitope specificity not captured by SSV shape features, (2) measurement noise in small-sample prototype estimation ($n = 5$), or (3) multimodal binding within the positive set. The size-matched control (E7b) confirms that these preferences are not artifacts of size confounding, supporting their biological validity even if predictive power is limited.

\textbf{Case 3: AB\_A1-245---Non-responder failure mode.} This antibody (8 positive glycans) exemplifies the non-responder category identified in E7c: no significant SSV dimension preferences detected despite adequate sample size, and very poor ranking performance (MRR = 0.058, Recall@5 = 0.0, mean rank = 34.8). Non-responder diagnostics (Table~\ref{tab:e7_diagnostics}) reveal high within-positive variance (mean SSV variance = 303.9 vs responders' mean = 14.9, Mann-Whitney $p = 0.016$). This indicates \emph{binding heterogeneity}: AB\_A1-245 recognizes structurally diverse glycans that do not share a consistent geometric profile. SSV features, which encode global shape descriptors, cannot discriminate such heterogeneous binding. This failure mode is biologically meaningful---it distinguishes \emph{narrow-specificity} agents (detectable via SSV) from \emph{broad-specificity} or \emph{sequence-motif-dependent} agents (requiring finer-grained representations). For AB\_A1-245, sequence-based methods (E8-ext) may outperform SSV: indeed, WURCS TF-IDF achieves MRR = 0.784 overall, suggesting that many agents bind via sequence epitopes rather than geometric fit.

\textbf{Interpretation.} These case studies illustrate three distinct regimes: (1) SSV-detectable shape preferences with strong predictive power (Case 1), (2) SSV-detectable preferences with limited prediction due to insufficient feature richness (Case 2), and (3) SSV-undetectable binding due to heterogeneity (Case 3). Together, they validate the benchmark's diagnostic utility: SSV features successfully identify \emph{when} geometric shape matters, even if they do not always achieve maximal predictive accuracy. The multi-prototype extension (E11, Sec.~\ref{sec:e11}) addresses Case 2-type limitations by allowing multiple structural modes per agent.

\textbf{Limitations.} All three cases rely on computationally generated structures (single-conformer PDB files from targeted acquisition); experimental 3D structures would strengthen interpretation. Epitope annotations from SugarBind are coarse-grained (e.g., ``T Antigen''), precluding fine-grained validation of SSV preferences against known binding sites. Future work should integrate structural binding data (PDB complexes) to validate SSV-preference concordance.
